% Orbit-Transfer — version française (traduction personnelle de paper.tex)
% Compilation: pdflatex paper-fr.tex (deux passages suffisent).
\documentclass[conference]{IEEEtran}
\IEEEoverridecommandlockouts
\usepackage{cite}
\usepackage[T1]{fontenc}
\usepackage{amsmath,amssymb,amsfonts}
\usepackage{textcomp}
\usepackage{booktabs}
\usepackage{url}
\usepackage{xcolor}
\usepackage[french]{babel}
\frenchbsetup{StandardLists=true}

% Étiquettes IEEEtran en français (résumé et mots-clés)
\providecommand{\IEEEabstractname}{Abstract}
\providecommand{\IEEEkeywordsname}{Index Terms}
\renewcommand{\IEEEabstractname}{Résumé}
\renewcommand{\IEEEkeywordsname}{Mots-clés}

\hyphenation{op-tical net-works semi-conduc-tor}

\begin{document}

\title{Orbit-Transfer : agrégation de bande passante\\ multi-relais avec codes fontaines rateless}

\author{\IEEEauthorblockN{Dylan Ondo}
\IEEEauthorblockA{\textit{Chercheur indépendant} \\
Libreville, Gabon \\
Email : danielebang59@gmail.com}}

\maketitle

\begin{abstract}
Ce papier présente Orbit-Transfer, un système de transfert de fichiers hybride
pair-à-pair et relais périphériques qui agrège la capacité de plusieurs chemins
réseau contraints à l'aide de codes fontaines rateless. Une charge utile est
découpée en $K$ symboles source, augmentés de $S$ contrôles de précodage LDPC
et de $H$ contrôles HDPC de type RaptorQ, et l'émetteur produit un flux
illimité de symboles encodés ; le
récepteur reconstruit la charge utile à partir de n'importe quels
$K+\varepsilon$ symboles distincts, sans demande de retransmission. Comme les
symboles de fontaine sont interchangeables et indépendants de l'ordre,
l'émetteur peut les répartir en tourniquet sur $N$ relais---et un canal
pair-à-pair direct optionnel---si bien que le débit agrégé est la somme des
chemins actifs. Nous évaluons la conception sur un banc d'essai multi-relais où
chaque relais est limité en débit par un régulateur à seau de jetons simulant un
lien montant contraint. Sur cinq essais, deux relais délivrent $2{,}2\times$ et
trois relais $3{,}2\times$ le débit d'un relais unique, avec un surcoût de
décodage négligeable dans la phase systématique sans perte. Nous rapportons
aussi un constat d'ingénierie pratique : les limiteurs de débit à attente
active sur le processeur s'effondrent sous la concurrence, et nous montrons
comment une minuterie haute résolution combinée à un seau de jetons
auto-correctif maintient des relais indépendants parfaitement uniformes.
\end{abstract}

\begin{IEEEkeywords}
codes d'effacement, codes fontaines, codes LT, codes Raptor, agrégation de
bande passante, transport multi-chemins, pair-à-pair, limitation de débit.
\end{IEEEkeywords}

\section{Introduction}
\IEEEPARstart{L}{es} transferts de grande taille sur des réseaux contraints se
heurtent à deux obstacles classiques : le lien goulot d'étranglement entre les
extrémités et les protocoles de retransmission qui multiplient latence et
gigue. Lorsqu'un chemin unique est limité en débit---un lien montant
résidentiel, un quota cellulaire ou un relais partagé---le remède évident est
d'en utiliser plusieurs à la fois. Le transport multi-chemins
(\texttt{MPTCP}~\cite{rfc6824}, SCTP) et le regroupement de
ressources~\cite{wischik2008survey} rendent cela possible au niveau de la
couche transport, mais ils exigent une coordination des chemins, un couplage de
congestion et un état de transport de bout en bout. En revanche, les codes
fontaines~\cite{byers1998digital,luby2002lt,mackay2005fountain} rendent
l'agrégation presque triviale : les symboles sont interchangeables, donc un
émetteur peut simplement les répartir en tourniquet sur chaque chemin
disponible et laisser le récepteur assembler la charge utile à partir de ce qui
arrive.

Orbit-Transfer exploite exactement cette propriété. L'émetteur découpe un
fichier en $K$ blocs source, applique un précodage de type Raptor---LDPC
puis HDPC de type RaptorQ~\cite{shokrollahi2006raptor,rfc6330}---pour obtenir
$L=K+S+H$ entrées précodées, et
émet un flux \emph{illimité} de symboles distincts. Le récepteur n'a besoin que
de $K+\varepsilon$ symboles \emph{quelconques} pour reconstruire le fichier ;
il ne demande jamais un bloc manquant. L'émetteur continue d'émettre jusqu'à ce
que le récepteur signale \textsc{Ready}, de sorte que toute perte, tout
réordonnancement ou toute défaillance de chemin est absorbé automatiquement.

La conséquence pratique est un ordonnanceur multi-chemins très simple :
répartir les symboles en tourniquet sur $N$ relais périphériques (et
optionnellement un canal pair-à-pair direct sur TCP ou QUIC), désactiver les
chemins en
panne, et compter sur la propriété rateless pour terminer. Comme les relais
n'ont pas besoin de se coordonner, la bande passante s'agrège linéairement :
$N$ relais de capacité $C$ délivrent jusqu'à $N\cdot C$, à condition que chaque
relais soit cadencé indépendamment et avec précision.

Ce papier apporte les contributions suivantes :
\begin{itemize}
\item Une description complète du système Orbit-Transfer : un protocole binaire
compact, un relais périphérique à salons avec mise en tampon pré-récepteur et
contre-pression, un ordonnanceur multi-chemins, et un chiffrement AEAD par
symbole (\S\ref{sec:arch}).
\item La conception du codec fontaine---phase systématique, distribution
Robust Soliton, précodage LDPC et HDPC (type Raptor/RaptorQ) et décodage par
épluchage---ainsi que ses paramètres (\S\ref{sec:codec}).
\item Une étude empirique de l'agrégation : deux relais délivrent $2{,}2\times$
et trois relais $3{,}2\times$ le débit d'un relais unique limité, et le codec
termine avec un surcoût quasi nul dans la phase systématique
(\S\ref{sec:eval}).
\item Une correction pratique de la limitation de débit : les seaux de jetons
à attente active ne passent pas à l'échelle sous concurrence ; les remplacer
par une minuterie haute résolution et un modèle de crédit auto-correctif
maintient des relais indépendants parfaitement uniformes
(\S\ref{sec:throttle}).
\end{itemize}

\section{Architecture du système}
\label{sec:arch}

\begin{figure}[!t]
\centering
\scriptsize
\begin{minipage}{0.95\columnwidth}
\begin{verbatim}
Emetteur --s--> Relais 1 --s--> Récepteur
  |                ...                 ^
  +--s--> Relais N --s-----------------+
  +--s--> direct (optionnel)-----------+
  <---------- READY -------------------+
\end{verbatim}
\end{minipage}
\caption{Modèle du système. L'émetteur répartit les symboles sur $N$ relais et
un canal pair-à-pair direct optionnel ; le récepteur fusionne tous les flux et
signale \textsc{Ready} une fois le décodage terminé.}
\label{fig:system}
\end{figure}

Le système (Fig.~\ref{fig:system}) comprend trois rôles : un \emph{émetteur},
un ou plusieurs \emph{relais} et un \emph{récepteur}. L'émetteur se connecte à
chaque relais via WebSocket et, lorsqu'un chemin direct est annoncé, au
récepteur via TCP brut ou QUIC (sélectionnés d'après le préfixe
\texttt{quic://} de l'adresse annoncée). Les deux parties partagent un
identifiant de session sur 64 bits.

\subsection{Protocole}
Les trames binaires transitent par WebSocket (liens relais) ou TCP brut (P2P
direct), avec l'agencement
\texttt{[type : u8][longueur : u32 LE][charge utile]}. Le
Tableau~\ref{tab:protocol} liste les types de messages. La trame \textsc{Meta}
porte le nom du fichier, la taille de la charge utile, la taille des symboles,
$K$ et une somme de contrôle XXH3 ; \textsc{Symbol} porte un symbole encodé ;
\textsc{Ready} termine le transfert ; \textsc{Direct} annonce l'adresse
d'écoute pair-à-pair du récepteur ; \textsc{Done} démonte une session.

\begin{table}[!t]
\centering
\caption{Messages du protocole.}
\label{tab:protocol}
\begin{tabular}{@{}lp{0.62\columnwidth}@{}}
\toprule
Message & Charge utile \\
\midrule
\textsc{Hello} & identifiant de session (u64), rôle (émetteur/récepteur) \\
\textsc{Meta} & identifiant, nom du fichier, taille, taille des symboles, $K$, somme \\
\textsc{Symbol} & identifiant, esi (u32), données du symbole \\
\textsc{Ready} & identifiant de session \\
\textsc{Direct} & identifiant, adresse annoncée \\
\textsc{Done} & identifiant de session \\
\bottomrule
\end{tabular}
\end{table}

\subsection{Relais}
Chaque relais maintient une table de \emph{salons} indexée par identifiant de
session, avec un lien émetteur et un lien récepteur par salon. Comme le
récepteur peut rejoindre plus tard, le relais met les symboles en tampon dans
une file bornée ($\le 256$\,Mio) et les vidange dans l'ordre à l'enregistrement.
Le routage repose uniquement sur l'identifiant de session ; le relais est
sans état vis-à-vis de la charge utile et ne la décode jamais. Les canaux de
sortie par connexion fournissent une contre-pression qui se propage au contrôle
de flux TCP, de sorte qu'un récepteur lent ou un lien limité ne provoque jamais
de pertes silencieuses---les symboles s'accumulent et l'émetteur est
finalement freiné plutôt que de perdre des données.

Le relais expose aussi un limiteur de débit de sortie par connexion
(\S\ref{sec:throttle}) utilisé dans nos expériences pour émuler un lien montant
contraint.

\subsection{Sessions émetteur et récepteur}
La session émetteur se connecte à tous les relais, ouvre des files bornées par
chemin (1024 symboles) avec une tâche d'écriture chacune, émet \textsc{Meta},
puis diffuse les symboles. Chaque tâche d'écriture vide sa file
indépendamment au rythme propre du chemin, si bien que l'agrégat est la somme
des chemins actifs ; un envoi en échec marque le chemin comme mort et sa charge
est redistribuée automatiquement.

La session récepteur lance une tâche de lecture par relais qui réémet les
trames sur un canal unique, de sorte qu'une boucle de décodage à base de
\texttt{select} fusionne tous les chemins (et le canal direct optionnel) en un
flux ordonné. Le chemin direct passe par TCP brut par défaut et par
\textbf{QUIC} lorsque le récepteur annonce une adresse \texttt{quic://}
(l'émetteur choisit le transport d'après le préfixe) ; les deux portent le même
encadrement binaire. Après décodage, elle envoie \textsc{Ready} sur chaque lien
puis \emph{vidange} chaque socket jusqu'à EOF avant de fermer---fermer avec des
données non lues émettrait une réinitialisation TCP et purgerait le
\textsc{Ready} en vol.

\section{Codec fontaine rateless}
\label{sec:codec}

\subsection{Encodage}
La charge utile est découpée en $K$ symboles source de taille égale. Un
précodage LDPC de type Raptor~\cite{shokrollahi2006raptor,rfc5053} calcule
\begin{equation}
S = \max\!\bigl(2,\ \lfloor K/100\rfloor + 10\bigr)
\end{equation}
symboles de contrôle ; le contrôle $i$ est le XOR des symboles source dont
l'indice est congru à $i$ modulo $S$, combiné aux sources $i\bmod K$ et
$(i{+}1)\bmod K$. Une couche HDPC de type RaptorQ~\cite{rfc6330} ajoute ensuite
$H$ contrôles haute densité, où $H$ est le plus petit entier tel que
\begin{equation}
\binom{H}{\lfloor H/2\rfloor} \ge K + S,
\end{equation}
et le contrôle $i$ XOR-e un sous-ensemble pseudo-aléatoire dense (environ la
moitié) des $K+S$ entrées précodées, ensemencé de manière déterministe depuis
l'indice du contrôle. Les $L=K+S+H$ symboles précodés sont les entrées d'un
encodeur LT~\cite{luby2002lt} qui tire son degré de la distribution Robust
Soliton avec $c=0{,}1$, $\delta=0{,}05$ et $R=c\sqrt{K}\ln(K/\delta)$. Les
contrôles HDPC servent de filet de sécurité au graphe de décodage : à faible
$K$, ils laissent la reconstruction se terminer près de $K+\varepsilon$ au lieu
de caler sur le graphe LDPC plus épars.

Le code est \emph{systématique} : pour esi $<K$, l'encodeur émet directement le
symbole source brut, si bien que dans le cas sans perte le récepteur
reconstruit la charge utile à partir des $K$ symboles systématiques avec un
surcoût \emph{nul}. Pour esi $\ge K$, il émet des combinaisons aléatoires. La
génération des symboles est déterministe et reproductible : la sélection des
voisins est ensemencée à partir de l'esi avec un PRNG ChaCha8, donc toute
partie peut recalculer le graphe d'encodage à partir d'un esi sans dictionnaire
partagé.

\subsection{Décodage}
Le décodeur construit le graphe de décodage sur les $L$ entrées précodées. Il
amorce le processus d'épluchage avec les $S$ équations de contrôle LDPC et les
$H$ équations HDPC, et itère les
équations de degré un : résoudre un symbole le XOR-e hors de chaque équation
qui le mentionne, et une équation de degré un se résout immédiatement. L'étape
d'épluchage est indexée---un index inverse entrée-à-équation fait qu'un symbole
résolu ne touche que les équations qui le mentionnent ($O(\text{degré})$ par
symbole) au lieu de parcourir toutes les équations actives, ce qui maintient le
décodeur WASM du navigateur rapide à grand $K$. Les
équations HDPC denses se résolvent en dernier, garantissant que le graphe se
termine. La reconstruction se termine lorsque les $K+S+H$ entrées sont
résolues ; la charge
utile est ensuite vérifiée contre la somme de contrôle XXH3 portée par
\textsc{Meta}. Les doublons et les symboles hors ordre sont inoffensifs (ils
sont dédupliqués par esi), ce qui est la propriété qui rend l'agrégation
multi-chemins triviale. Le noyau XOR est vectorisable SIMD, et le débit
d'encodage dépasse $3$\,Gbit/s sur du matériel courant avec une charge utile de
64\,Mio.

\section{Agrégation multi-chemins}
\label{sec:agg}

L'agrégation dans Orbit-Transfer se réduit à une décision d'ordonnancement,
car les symboles de fontaine sont interchangeables et indépendants de l'ordre.
L'émetteur utilise un ordonnanceur déterministe~\cite{maymounkov2002online} qui
répartit la charge des relais en tourniquet sur les $N$ relais tout en
conservant un rapport relais/direct. Deux propriétés comptent en pratique :
\begin{itemize}
\item \textbf{Bascule en cas de chemin mort.} Un relais qui meurt en cours de
transfert est désactivé ; ses symboles en file sont perdus, mais le récepteur
ne s'en aperçoit jamais---il a besoin de $K+\varepsilon$ symboles
quelconques, pas d'un sous-ensemble précis. L'ordonnanceur passe simplement à
la suite.
\item \textbf{Arrêt par vidange.} Après \textsc{Ready}, le récepteur vidange
chaque lien relais jusqu'à EOF afin que les relais purgent leurs symboles en
file proprement plutôt que de subir une réinitialisation TCP.
\end{itemize}

Comme chaque relais est cadencé indépendamment, l'agrégat attendu est $N\cdot
C$. La section expérimentale (\S\ref{sec:eval}) confirme que cela tient à la
variance près des limiteurs individuels.

\section{Limitation de débit des relais}
\label{sec:throttle}

Nos expériences émulent des liens montants contraints avec un régulateur à seau
de jetons par connexion au niveau du relais : chaque symbole coûte son nombre
d'octets et le seau se remplit au débit configuré. Deux détails
d'implémentation se sont révélés critiques.

\subsection{L'attente active ne passe pas à l'échelle}
La première implémentation cadencait un symbole en attendant activement sur
\texttt{Instant::now()} tout en cédant la main à l'ordonnanceur :
\begin{verbatim}
while Instant::now() < deadline {
    yield_now().await;
}
\end{verbatim}
Seul, un relais était cadencé à $\approx 343$ symboles/s au lieu des $500$/s
nominaux. Sous trois relais concurrents, les débits par relais se sont
effondrés à $361/316/350$ symboles/s ($1028$/s au total au lieu de $1500$/s) :
chaque tâche d'attente se dispute le processeur et le surcoût d'ordonnancement
croît avec le nombre de boucles concurrentes. Le limiteur sous-estime le débit
et perd son indépendance, ce qui sape directement l'affirmation d'agrégation.

\subsection{Sommeils parkés avec un seau auto-correctif}
Nous avons remplacé l'attente active par un vrai \texttt{sleep} et porté la
résolution de la minuterie Windows à 1\,ms via \texttt{timeBeginPeriod(1)}.
Dormir parke la tâche, donc les relais concurrents ne se disputent plus le
processeur. Pourtant, la minuterie \texttt{tokio} dort encore $\approx 3{,}7$\,ms
pour un $2$\,ms demandé, ce qui à lui seul sous-estimerait le débit nominal de
moitié.

La correction est un seau de jetons dont la comptabilité est
\emph{auto-corrective} : la fenêtre de crédit est close au \emph{début} de
chaque appel et \emph{jamais} après le sommeil, si bien que le temps réellement
pris par le sommeil est recrédité au symbole suivant. Le débit à long terme
converge donc vers la valeur configurée quelle que soit la granularité de la
minuterie. Un plafond de rafale ($\approx 10$\,ms de crédit) empêche une longue
inactivité de débloquer une rafale illimitée. Le résultat est un débit uniforme
de $613/613/614$ symboles/s sur trois relais concurrents---un passage à
l'échelle parfait de $3\times$ sans contention.

\section{Évaluation expérimentale}
\label{sec:eval}

\subsection{Configuration}
Toutes les expériences s'exécutent sur une seule machine Windows 11 en
loopback, en mode release (les builds de débogage ralentissent l'encodeur d'un
ordre de grandeur). Chaque relais limite son trafic sortant à $2$\,Mio/s
nominaux. La charge utile est de 16\,Mio avec des symboles de 4096 octets
($K=4096$). Nous rapportons cinq essais du banc d'agrégation ; le tableau donne
la moyenne et l'étendue. Les facteurs d'agrégation sont calculés par essai
comme le rapport des temps écoulés (la taille de la charge utile est identique
entre les cas).

\subsection{Surcoût du codec}
Dans le cas sans perte, la phase systématique reconstruit à elle seule la
charge utile : le décodeur termine exactement à $K$ symboles (surcoût $1{,}000$).
Sous perte simulée de $5\%$, le surcoût est d'environ $1{,}1\times L$, bien dans
la garantie Raptor. Le débit d'encodage dépasse $3$\,Gbit/s pour une charge
utile de 64\,Mio ; le décodage de la même charge utile se termine en environ
$0{,}13$\,s.

\subsection{Agrégation}
\begin{table}[!t]
\centering
\caption{Débit en fonction du nombre de relais limités (charge utile de
16\,Mio, symboles de 4096 octets, moyenne de 5 essais).}
\label{tab:bench}
\begin{tabular}{@{}lcccc@{}}
\toprule
Cas & Mio/s & Étendue & Symboles & Facteur \\
\midrule
1 relais & 2.29 & 2.19--2.34 & $\approx$21.6k & 1.00 \\
2 relais & 5.01 & 4.31--5.39 & $\approx$34.5k & 2.19 \\
3 relais & 7.24 & 6.66--7.78 & $\approx$34.6k & 3.16 \\
3 relais + direct & 5.76 & 5.40--6.30 & $\approx$43.9k & 2.51 \\
\bottomrule
\end{tabular}
\end{table}

Le Tableau~\ref{tab:bench} rapporte les résultats. Deux relais délivrent
$2{,}19\times$ (étendue $1{,}84$--$2{,}47\times$) et trois relais $3{,}16\times$
(étendue $2{,}84$--$3{,}70\times$) le débit de référence à un seul relais, ce
qui correspond à l'échelle attendue $N\times$ à la variance des limiteurs près.
Le cas à un relais est remarquablement stable ($2{,}19$--$2{,}34$\,Mio/s sur
les essais).

Le cas « 3 relais + direct » ajoute l'écouteur pair-à-pair du récepteur ; le
canal direct ne transporte que la queue redondante (esi $\ge K$), car le
récepteur se complète à environ $K$ symboles et les relais limités saturent
déjà le flux de décodage. Par conséquent, le P2P direct ajoute peu sur ce banc
de test en loopback. Il est prévu pour le scénario opposé---quand les relais
sont indisponibles ou constituent le goulot---comme voie rapide avec bascule
automatique, plutôt que comme multiplicateur par-dessus des relais saturés.

\subsection{Indépendance des limiteurs}
Le Tableau~\ref{tab:throttle} isole le limiteur en cadencant trois relais en
parallèle sans sessions ni décodage. Les débits par relais sont à $0{,}4\%$
près les uns des autres et leur somme vaut $3\times$ le débit d'un relais
unique, ce qui confirme que l'agrégation n'est limitée que par les conduits
individuels et non par une ressource processeur partagée.

\begin{table}[!t]
\centering
\caption{Trafic sortant par relais avec trois relais limités simultanément.}
\label{tab:throttle}
\begin{tabular}{@{}lcc@{}}
\toprule
Relais & Débit par relais & Total \\
\midrule
1 & 624 sym/s & 624/s \\
3 & 613, 613, 614 sym/s & 1841/s \\
\bottomrule
\end{tabular}
\end{table}

\section{Travaux connexes}
Les codes fontaines naissent avec le cadre de la fontaine numérique de
Byers \emph{et al.}~\cite{byers1998digital} et la construction LT de
Luby~\cite{luby2002lt} ; MacKay en donne un panorama
accessible~\cite{mackay2005fountain}. Les codes Raptor~\cite{shokrollahi2006raptor,rfc5053}
et RaptorQ~\cite{rfc6330} ajoutent un précodage fixe pour ramener le surcoût de
décodage à quelques pour cent et sont utilisés dans les normes de diffusion
3GPP et de livraison fiable. Notre codec suit la recette Raptor (LT sur un
précodage LDPC complété d'une couche HDPC de type RaptorQ, phase systématique)
à plus petite échelle et avec une graine de
voisins déterministe et reproductible.

Le transport multi-chemins a été normalisé sous la forme de
\texttt{MPTCP}~\cite{rfc6824} et étudié à travers le prisme du regroupement de
ressources de Wischik \emph{et al.}~\cite{wischik2008survey}. Contrairement au
couplage au niveau transport, notre schéma n'a besoin d'aucune coordination de
congestion entre les chemins : le code rateless absorbe les différences de
débit entre chemins, et le seul état partagé est l'identifiant de session.

Les codes fontaines sont proposés depuis longtemps pour la diffusion de masse
pair-à-pair~\cite{maymounkov2002online} ; Orbit-Transfer rend l'extension
multi-chemins explicite et la quantifie sur des relais limités en débit, ce qui,
à notre connaissance, n'est pas rapporté dans la littérature des fontaines.

\section{Conclusion et travaux futurs}
Orbit-Transfer montre que les codes fontaines rateless transforment
l'agrégation multi-chemins en un problème d'ordonnancement au gain
essentiellement linéaire : $N$ relais limités en débit délivrent environ
$N\times$ le débit d'un seul, car les symboles sont interchangeables et chaque
chemin est cadencé indépendamment. Les enseignements d'ingénierie sont (i) le
précodage systématique de type Raptor supprime entièrement le surcoût dans le
cas sans perte, et (ii) les limiteurs de débit doivent reposer sur des sommeils
parkés et être auto-correctifs pour maintenir des chemins indépendants
uniformes sous concurrence.

Les travaux en cours et futurs incluent l'extension du chemin direct au-delà de
TCP et QUIC (déjà implémenté : le récepteur peut annoncer un écouteur
\texttt{quic://} et l'émetteur choisit QUIC automatiquement, ce qui offre une
connexion 0-RTT et la migration de connexion), et la poursuite du précodage de
type RaptorQ vers un surcoût encore plus faible à très petit $K$. Le relais est
déjà agnostique au transport, et le protocole est assez compact pour transiter
sur tout canal ordonné ou non.

En complément de la pile multi-chemins native, un pont WebRTC est livré avec le
paquet WebAssembly : les transferts navigateur-à-navigateur s'exécutent sur un
canal de données SCTP avec le même codec fontaine compilé en WASM. La
signalisation est automatisée via un minuscule salon en mémoire sur le serveur
statique : l'émetteur glisse un fichier et obtient un lien partageable
\texttt{?room=<id>}, le récepteur l'ouvre, et l'offre/réponse SDP transite par
HTTP. Une leçon pratique est ressortie
des tests navigateur : Chrome masque les candidats ICE hôtes derrière des noms
mDNS (\texttt{*.local}) que certaines interfaces réseau ne peuvent pas
résoudre, laissant l'agent ICE sans paire de candidats (bloqué dans l'état
« new ») ; comme les ports des candidats ne sont pas masqués, l'ajout explicite
d'un candidat loopback pointant vers le port du pair rétablit la connectivité
entre onglets de la même machine. Le pont est validé de bout en bout dans le
navigateur et par des tests d'intégration sur canal simulé (y compris le flux
automatique par lien partageable). La version actuelle repose sur des candidats
ICE hôtes uniquement, si bien que le flux automatisé est validé sur une seule
machine (loopback) et sur un LAN ; le pair-à-pair sur Internet public exige
STUN/TURN, ce qui---avec un transport par datagrammes QUIC/WebTransport comme
alternative au canal SCTP offrant un débit supérieur, et un basculement vers le
relais quand la connectivité directe échoue---fait l'objet de travaux en cours.

\begin{thebibliography}{1}
\bibitem{byers1998digital}
J.~W. Byers, M.~Luby, M.~Mitzenmacher, and A.~Rege, ``A digital fountain approach
to reliable distribution of bulk data,'' in \emph{Proc. ACM SIGCOMM}, 1998,
pp.~56--67.
\bibitem{luby2002lt}
M.~Luby, ``LT codes,'' in \emph{Proc. 43rd IEEE Symp. Foundations of Computer
Science (FOCS)}, 2002, pp.~271--280.
\bibitem{mackay2005fountain}
D.~J.~C. MacKay, ``Fountain codes,'' \emph{IEE Proc. Communications},
vol.~152, no.~6, pp.~1062--1068, 2005.
\bibitem{shokrollahi2006raptor}
A.~Shokrollahi, ``Raptor codes,'' \emph{IEEE Trans. Information Theory},
vol.~52, no.~6, pp.~2551--2567, 2006.
\bibitem{rfc5053}
M.~Luby, A.~Shokrollahi, M.~Watson, and T.~Stockhammer, ``Raptor forward error
correction scheme for object delivery,'' IETF RFC 5053, 2007.
\bibitem{rfc6330}
M.~Luby, A.~Shokrollahi, M.~Watson, and T.~Stockhammer, ``RaptorQ forward error
correction scheme for object delivery,'' IETF RFC 6330, 2011.
\bibitem{rfc6824}
A.~Ford, C.~Raiciu, M.~Handley, and O.~Bonaventure, ``TCP extensions for
multipath operation with multiple addresses,'' IETF RFC 6824, 2013.
\bibitem{wischik2008survey}
D.~Wischik, M.~Handley, and M.~Bagnulo Braun, ``The resource pooling
principle,'' \emph{ACM SIGCOMM Comput. Commun. Rev.}, vol.~38, no.~5,
pp.~47--52, 2008.
\bibitem{maymounkov2002online}
P.~Maymounkov and D.~Mazi{\`e}res, ``Rateless codes and big downloads,'' in
\emph{Proc. 2nd Int. Workshop on Peer-to-Peer Systems (IPTPS)}, 2003,
pp.~247--255.
\bibitem{webrtc}
A.~Bergkvist \emph{et al.}, ``WebRTC 1.0: real-time communication between
browsers,'' W3C Recommendation, 2021.
\end{thebibliography}

\end{document}